\documentclass[11pt]{article}
\usepackage{docstyle}
% 12 mm side margins: exam questions print at natural size (crops are at most 181 mm wide)
\newgeometry{left=12mm,right=12mm,top=14mm,bottom=20mm,footskip=9mm}
\begin{document}
\sheethead{Lesson 5 \textperiodcentered\ Homework}{NCEA Level 2 \textperiodcentered\ AS 91262 \textperiodcentered\ Mixed revision and Excellence questions}{about 60 minutes}{}
Show all working. Justify every maximum or minimum; give units and meaning for rates. [A] Achieved, [M] Merit, [E] Excellence.

\Q Differentiate $y=(x+2)(x^{2}-3)$.\mk{A}
\ALine{$y=x^{3}+2x^{2}-3x-6$}
\ALineT{$\dydx=3x^{2}+4x-3$}

\Q Find the \tturns of $y=x^{3}-6x^{2}+9$ and justify their nature.\mk{A}
\ALineT{$\dydx=3x^{2}-12x=3x(x-4)=0$, so $x=0$ or $x=4$}
\ALine{$y(0)=9$;\quad $y(4)=64-96+9=-23$}
\ALineT{$\dfrac{\mathrm{d}^{2}y}{\mathrm{d}x^{2}}=6x-12$:\ at $x=0$ it is $-12<0$, so $(0,\,9)$ is a maximum}
\ALine{at $x=4$ it is $12>0$, so $(4,\,-23)$ is a minimum}

\Q $f'(x)=6x-2$ and the curve passes through $(-1,\,8)$. Find $f(x)$.\mk{A}
\ALine{$f(x)=3x^{2}-2x+c$;\ \ $3+2+c=8$, so $c=3$}
\ALine{$f(x)=3x^{2}-2x+3$}

\Q $f'(x)=4-x^{2}$. Without drawing, describe the \tturns of $f$: where they are, and which is a maximum.\mk{A}
\ALine{$f'(x)=0$ when $x=2$ or $x=-2$}
\ALine{At $x=-2$, $f'$ changes from $-$ to $+$: minimum}
\ALine{At $x=2$, $f'$ changes from $+$ to $-$: maximum}

\Q The \ttangent to $y=x^{2}+bx+c$ at the point $(1,\,3)$ has gradient $5$. Find $b$ and $c$.\mk{M}
\ALineT{Gradient: $\dydx=2x+b$;\ at $x=1$:\ $2+b=5$, so $b=3$}
\ALine{Point: $1+3+c=3$, so $c=-1$}

\Q A particle has acceleration $a=4-2t\mpss$, starts at the origin, and has velocity $5\mps$ when $t=0$. Find its maximum velocity and its displacement at that moment.\mk{M}
\ALine{$v=4t-t^{2}+5$ (using $v(0)=5$)}
\ALine{Maximum velocity when $a=0$: $t=2$ s;\ \ $v(2)=8-4+5=9\mps$ (maximum: $a$ changes from $+$ to $-$)}
\ALineT{$s=2t^{2}-\dfrac{1}{3}t^{3}+5t$ (from the origin);\ \ $s(2)=8-\dfrac{8}{3}+10=15.3$ m}

\ppQ{[E]}
\ppq{a21_1c3}

\ppQ{[E]}
\ppq{a22_3c}

\clearpage
{\sffamily\bfseries Question \theqn\ continued}\quad working space
\medskip
\ALine{$S'(b)=2\sqrt{3}(30-2b)(-2)+4\sqrt{3}\,b=\sqrt{3}(12b-120)$}
\ALine{$S'(b)=0$ when $b=10$, so $a=30-2(10)=10$}
\ALine{$S''(b)=12\sqrt{3}>0$, so the total surface area is a minimum}
\ALine{$a=10$ cm and $b=10$ cm (minimum total surface area $300\sqrt{3}=519.6\ \mathrm{cm}^{2}$)}
\ALines{14}

\ppQ{Extension [E]}
\ppq{a21_3d2}
\end{document}
